% Options for packages loaded elsewhere
% Options for packages loaded elsewhere
\PassOptionsToPackage{unicode}{hyperref}
\PassOptionsToPackage{hyphens}{url}
\PassOptionsToPackage{dvipsnames,svgnames,x11names}{xcolor}
%
\documentclass[
  11pt,
  a4paper,
  DIV=11,
  numbers=noendperiod]{scrartcl}
\usepackage{xcolor}
\usepackage[lmargin=2cm,rmargin=2cm,tmargin=2cm,bmargin=2cm]{geometry}
\usepackage{amsmath,amssymb}
\setcounter{secnumdepth}{-\maxdimen} % remove section numbering
\usepackage{iftex}
\ifPDFTeX
  \usepackage[T1]{fontenc}
  \usepackage[utf8]{inputenc}
  \usepackage{textcomp} % provide euro and other symbols
\else % if luatex or xetex
  \usepackage{unicode-math} % this also loads fontspec
  \defaultfontfeatures{Scale=MatchLowercase}
  \defaultfontfeatures[\rmfamily]{Ligatures=TeX,Scale=1}
\fi
\usepackage{lmodern}
\ifPDFTeX\else
  % xetex/luatex font selection
\fi
% Use upquote if available, for straight quotes in verbatim environments
\IfFileExists{upquote.sty}{\usepackage{upquote}}{}
\IfFileExists{microtype.sty}{% use microtype if available
  \usepackage[]{microtype}
  \UseMicrotypeSet[protrusion]{basicmath} % disable protrusion for tt fonts
}{}
\makeatletter
\@ifundefined{KOMAClassName}{% if non-KOMA class
  \IfFileExists{parskip.sty}{%
    \usepackage{parskip}
  }{% else
    \setlength{\parindent}{0pt}
    \setlength{\parskip}{6pt plus 2pt minus 1pt}}
}{% if KOMA class
  \KOMAoptions{parskip=half}}
\makeatother
% Make \paragraph and \subparagraph free-standing
\makeatletter
\ifx\paragraph\undefined\else
  \let\oldparagraph\paragraph
  \renewcommand{\paragraph}{
    \@ifstar
      \xxxParagraphStar
      \xxxParagraphNoStar
  }
  \newcommand{\xxxParagraphStar}[1]{\oldparagraph*{#1}\mbox{}}
  \newcommand{\xxxParagraphNoStar}[1]{\oldparagraph{#1}\mbox{}}
\fi
\ifx\subparagraph\undefined\else
  \let\oldsubparagraph\subparagraph
  \renewcommand{\subparagraph}{
    \@ifstar
      \xxxSubParagraphStar
      \xxxSubParagraphNoStar
  }
  \newcommand{\xxxSubParagraphStar}[1]{\oldsubparagraph*{#1}\mbox{}}
  \newcommand{\xxxSubParagraphNoStar}[1]{\oldsubparagraph{#1}\mbox{}}
\fi
\makeatother


\usepackage{longtable,booktabs,array}
\newcounter{none} % for unnumbered tables
\usepackage{calc} % for calculating minipage widths
% Correct order of tables after \paragraph or \subparagraph
\usepackage{etoolbox}
\makeatletter
\patchcmd\longtable{\par}{\if@noskipsec\mbox{}\fi\par}{}{}
\makeatother
% Allow footnotes in longtable head/foot
\IfFileExists{footnotehyper.sty}{\usepackage{footnotehyper}}{\usepackage{footnote}}
\makesavenoteenv{longtable}
\usepackage{graphicx}
\makeatletter
\newsavebox\pandoc@box
\newcommand*\pandocbounded[1]{% scales image to fit in text height/width
  \sbox\pandoc@box{#1}%
  \Gscale@div\@tempa{\textheight}{\dimexpr\ht\pandoc@box+\dp\pandoc@box\relax}%
  \Gscale@div\@tempb{\linewidth}{\wd\pandoc@box}%
  \ifdim\@tempb\p@<\@tempa\p@\let\@tempa\@tempb\fi% select the smaller of both
  \ifdim\@tempa\p@<\p@\scalebox{\@tempa}{\usebox\pandoc@box}%
  \else\usebox{\pandoc@box}%
  \fi%
}
% Set default figure placement to htbp
\def\fps@figure{htbp}
\makeatother





\setlength{\emergencystretch}{3em} % prevent overfull lines

\providecommand{\tightlist}{%
  \setlength{\itemsep}{0pt}\setlength{\parskip}{0pt}}



 


\KOMAoption{captions}{tableheading}
\makeatletter
\@ifpackageloaded{caption}{}{\usepackage{caption}}
\AtBeginDocument{%
\ifdefined\contentsname
  \renewcommand*\contentsname{Table of contents}
\else
  \newcommand\contentsname{Table of contents}
\fi
\ifdefined\listfigurename
  \renewcommand*\listfigurename{List of Figures}
\else
  \newcommand\listfigurename{List of Figures}
\fi
\ifdefined\listtablename
  \renewcommand*\listtablename{List of Tables}
\else
  \newcommand\listtablename{List of Tables}
\fi
\ifdefined\figurename
  \renewcommand*\figurename{Figure}
\else
  \newcommand\figurename{Figure}
\fi
\ifdefined\tablename
  \renewcommand*\tablename{Table}
\else
  \newcommand\tablename{Table}
\fi
}
\@ifpackageloaded{float}{}{\usepackage{float}}
\floatstyle{ruled}
\@ifundefined{c@chapter}{\newfloat{codelisting}{h}{lop}}{\newfloat{codelisting}{h}{lop}[chapter]}
\floatname{codelisting}{Listing}
\newcommand*\listoflistings{\listof{codelisting}{List of Listings}}
\makeatother
\makeatletter
\makeatother
\makeatletter
\@ifpackageloaded{caption}{}{\usepackage{caption}}
\@ifpackageloaded{subcaption}{}{\usepackage{subcaption}}
\makeatother
\usepackage{bookmark}
\IfFileExists{xurl.sty}{\usepackage{xurl}}{} % add URL line breaks if available
\urlstyle{same}
\hypersetup{
  pdftitle={Project Göç Dinamikleri},
  colorlinks=true,
  linkcolor={blue},
  filecolor={Maroon},
  citecolor={Blue},
  urlcolor={Blue},
  pdfcreator={LaTeX via pandoc}}


\title{Project Göç Dinamikleri}
\author{}
\date{}
\begin{document}
\maketitle


Welcome to our project page.

Our Team:

Selin DÜZEN Burcu BARIŞ

Keep an eye on this space to stay updated with my project activities.

\section{1. Project Overview and
Scope}\label{project-overview-and-scope}

Göç, yalnızca bireylerin coğrafi yer değiştirmesi değil; ekonomik
arayışların, sosyal imkanların ve yapısal kalitenin bir bölgedeki
izdüşümüdür. Türkiye'de iller arasındaki gelişmişlik ve yaşam
standartları farkı, iç göç hareketlerinin en birincil dinamiklerinden
birini oluşturmaktadır. Bu projede, TÜİK verilerini kullanarak
Türkiye'nin 81 ilinin Net Göç (Net Göç = Göç Alan - Göç Veren) rakamları
ile İllerde Yaşam Endeksi (Yaşanabilirlik) göstergeleri arasındaki çok
boyutlu ilişkiyi analitik yöntemlerle incelemeyi hedefliyoruz. Amacımız,
insanların hangi yaşamsal faktörler (ekonomi, eğitim, sağlık, çevre
kalitesi, güvenlik vb.) nedeniyle mevcut şehirlerini terk ettiğini veya
belirli şehirlere yöneldiğini veri analitiği ve istatistiksel modellerle
ortaya koymaktır. Bu çalışma, bölgesel kalkınma politikaları ve kentsel
planlama kararları için veri odaklı bir projeksiyon sunması açısından
kritik bir karar destek mekanizması niteliğindedir. Projenin Temel
Amaçları Göç Profillerinin Çıkarılması: Türkiye'deki ``Net Göç
Şampiyonları'' (en çok net göç alan) ve ``Göç Kaynakları'' (en çok net
göç veren) illeri tespit etmek ve bunları coğrafi olarak
haritalandırmak. Yaşanabilirlik Boyutlarının Etkisi: Yaşam endeksini
oluşturan alt kırılımlardan (Gelir ve Servet, Sağlık, Eğitim, Çevre,
Güvenlik vb.) hangilerinin göç kararı üzerinde daha baskın birer
``itici'' veya ``çekici'' güç olduğunu istatistiksel olarak kanıtlamak.
Kümeleme (Clustering): İlleri hem göç davranışlarına hem de
yaşanabilirlik skorlarına göre homojen gruplara ayırarak benzer
sosyo-ekonomik dinamiklere sahip şehir yapılarını keşfetmek.

\section{2. Data}\label{data}

xxxxxx

\#veri seti, veri kümenizi açıklayın. Herhangi bir tam analiz yapmanız
veya sonuç çıkarmanız gerekmemekle birlikte, verilere kapsamlı bir genel
bakış sunmayı hedefleyin---sütunların ne anlama geldiğini, değerlerin
nasıl dağıldığını açıklayın ve potansiyel analizler hakkında
fikirlerinizi paylaşın. Özlü tutun, bazı grafikler ekleyebilirsiniz.

\begin{enumerate}
\def\labelenumi{\arabic{enumi}.}
\tightlist
\item
  İllerin yıllara sari popülasyon değeri
\item
  Göç etme nedenlerine göre illerin aldığı göç miktarı
\item
  Göç etme nedenlerine göre illerin verdiği göç miktarı
\item
  Göç etme nedenlerine göre illerin net göç miktarı
\end{enumerate}

\subsection{2.1 Data Source}\label{data-source}

xxxxxx

\subsection{2.2 General Information About
Data}\label{general-information-about-data}

xxxxxx

\subsection{2.3 Reason of Choice}\label{reason-of-choice}

xxxxxx

\subsection{2.4 Preprocessing}\label{preprocessing}

\#Farklı excel verilerinde yer alan İllerin aldığı ve Verdiği Göç
Değerlerinden Net Göç Hesaplanması ve verilerdeki char formatında
yazılan boşluklu değerlerin numaric değerler olarak güncellenmesi
işlemleri uygulanmıştır.

\#install.packages(``ggplot2'') \#install.packages(``dplyr'')\\
\#install.packages(``openxlsx'') library(ggplot2) library(dplyr)
library(readxl) library(writexl) library(stringr) library(forcats)
library(tidyr) library(openxlsx)

\section{Net göç bilgisi hesaplama ve excel dosyası
oluşturma}\label{net-guxf6uxe7-bilgisi-hesaplama-ve-excel-dosyasux131-oluux15fturma}

VerilenGoc \textless-
read\_excel(``C:/EMU660\_Project/emu660-spring2026-Selindznn/VerilenGoc.xlsx'',
col\_names = TRUE) AlinanGoc \textless-
read\_excel(``C:/EMU660\_Project/emu660-spring2026-Selindznn/AlinanGoc.xlsx'',
col\_names = TRUE) toplam\_sutun \textless- ncol(VerilenGoc) fark\_df
\textless- VerilenGoc for (i in 3:toplam\_sutun) \{ v\_clean \textless-
as.character(VerilenGoc{[}{[}i{]}{]}) v\_clean \textless-
str\_replace\_all(v\_clean, ``\textbackslash s'', ``\,``)
VerilenGoc{[}{[}i{]}{]} \textless- as.numeric(v\_clean)\\
a\_clean \textless- as.character(AlinanGoc{[}{[}i{]}{]}) a\_clean
\textless- str\_replace\_all(a\_clean,''\textbackslash s'', ``\,``)\\
AlinanGoc{[}{[}i{]}{]} \textless- as.numeric(a\_clean) \} fark\_df{[},
3:toplam\_sutun{]} \textless- VerilenGoc{[}, 3:toplam\_sutun{]} -
AlinanGoc{[}, 3:toplam\_sutun{]} for (i in 3:toplam\_sutun) \{
fark\_df{[}{[}i{]}{]} \textless- as.numeric(fark\_df{[}{[}i{]}{]}) \}
write\_xlsx(fark\_df,''C:/EMU660\_Project/emu660-spring2026-Selindznn/NetGoc.xlsx'')

print(VerilenGoc{[}, 3:toplam\_sutun{]})

\section{İllerin populasyon verisi ile Net göç hızı hesaplama ve excel
dosyası
oluşturma}\label{illerin-populasyon-verisi-ile-net-guxf6uxe7-hux131zux131-hesaplama-ve-excel-dosyasux131-oluux15fturma}

kaynak\_yolu \textless-
``C:/EMU660\_Project/emu660-spring2026-Selindznn/NetGoc.xlsx''
hedef\_yolu \textless-
``C:/EMU660\_Project/emu660-spring2026-Selindznn/IlNufuslarii.xlsx''

df\_kaynak \textless- read\_excel(kaynak\_yolu, col\_names = TRUE)

NetGocVerisi \textless- df\_kaynak{[}1:82, 3, drop = TRUE{]}

wb \textless- loadWorkbook(hedef\_yolu)

writeData(wb, sheet = 1, x = ``Net Göç'', startCol = 3, startRow = 1,
colNames = FALSE) writeData(wb, sheet = 1, x = NetGocVerisi, startCol =
3, startRow = 2, colNames = FALSE)

saveWorkbook(wb, hedef\_yolu, overwrite = TRUE)

\section{3. Analysis}\label{analysis}

\subsection{3.1 Exploratory Data
Analysis}\label{exploratory-data-analysis}

\section{Yıllara göre Toplam Verilen Göç Miktarı
Değişimi}\label{yux131llara-guxf6re-toplam-verilen-guxf6uxe7-miktarux131-deux11fiux15fimi}

grafik\_verisi \textless- VerilenGoc{[}VerilenGoc{[}, 2{]} ==
``Toplam-Total'', {]} print(grafik\_verisi) head(grafik\_verisi)

yil \textless- factor(grafik\_verisi{[}{[}1{]}{]}, levels =
sort(unique(grafik\_verisi{[}{[}1{]}{]}), decreasing = TRUE))
goc\_miktari \textless- as.numeric(grafik\_verisi{[}{[}3{]}{]})/1000

ggplot(data = grafik\_verisi, aes(x = yil, y = goc\_miktari)) +
geom\_col(fill = ``steelblue'', alpha = 0.7) + geom\_line(color =
``black'', size = 2) + geom\_point(color = ``darkred'', size = 3) +
geom\_text(aes(label = round(goc\_miktari, 1)), vjust = -0.5, size =
3.5) + labs( title = ``Yıllara Göre Toplam Verilen Göç Miktarı'', x =
``Yıl'', y = ``Göç Miktarı (x1.000)'' )

\section{Ankara için Yıllara Göre Toplam Verilen Göç
Miktarı}\label{ankara-iuxe7in-yux131llara-guxf6re-toplam-verilen-guxf6uxe7-miktarux131}

grafik\_verisi \textless- VerilenGoc{[}VerilenGoc{[}, 2{]} ==
``Ankara'', {]} print(grafik\_verisi) head(grafik\_verisi)

yil \textless- factor(grafik\_verisi{[}{[}1{]}{]}, levels =
sort(unique(grafik\_verisi{[}{[}1{]}{]}), decreasing = TRUE))
goc\_miktari \textless- as.numeric(grafik\_verisi{[}{[}3{]}{]})/1000

ggplot(data = grafik\_verisi, aes(x = yil, y = goc\_miktari)) +
geom\_col(fill = ``steelblue'', alpha = 0.7) + geom\_line(color =
``black'', size = 2) + geom\_point(color = ``darkred'', size = 3) +
geom\_text(aes(label = round(goc\_miktari, 1)), vjust = -0.5, size =
3.5) + labs( title = ``Ankara için Yıllara Göre Toplam Verilen Göç
Miktarı'', x = ``Yıl'', y = ``Göç Miktarı (x1.000)'' )

\section{2024 yılı İllerin Alınan Göç
Miktarları}\label{yux131lux131-illerin-alux131nan-guxf6uxe7-miktarlarux131}

AlinanGoc \textless-
read\_excel(``C:/EMU660\_Project/emu660-spring2026-Selindznn/AlinanGoc.xlsx'',col\_names
= TRUE)

colnames(AlinanGoc){[}1:3{]} \textless- c(``Yil'', ``Il'',
``Goc\_Miktari'') veri\_2024 \textless- subset(AlinanGoc, Yil == 2024 \&
Il != ``Toplam-Total'')
veri\_2024\(Goc_Miktari_Bin <- veri_2024\)Goc\_Miktari / 1000

ggplot(veri\_2024, aes(x = Goc\_Miktari\_Bin, y = reorder(Il,
Goc\_Miktari\_Bin))) + geom\_col(aes(fill = Goc\_Miktari\_Bin
\textgreater{} 0)) + scale\_fill\_manual(values = c(``coral'',
``skyblue''), guide = ``none'') + theme\_minimal() + labs( title =
``2024 Yılı İllerin Alınan Göç Miktarları'', x = ``Alınan Göç Miktarı
(Bin Kişi)'', y = ``İller'' )

\section{2024 yılı İllerin net Göç
Miktarları}\label{yux131lux131-illerin-net-guxf6uxe7-miktarlarux131}

NetGoc \textless-
read\_excel(``C:/EMU660\_Project/emu660-spring2026-Selindznn/NetGoc.xlsx'',col\_names
= TRUE)

colnames(NetGoc){[}1:3{]} \textless- c(``Yil'', ``Il'',
``Goc\_Miktari'') veri\_2024 \textless- subset(NetGoc, Yil == 2024 \& Il
!= ``Toplam-Total'')
veri\_2024\(Goc_Miktari_Bin <- veri_2024\)Goc\_Miktari / 1000

ggplot(veri\_2024, aes(x = Goc\_Miktari\_Bin, y = reorder(Il,
Goc\_Miktari\_Bin))) + geom\_col(aes(fill = Goc\_Miktari\_Bin
\textgreater{} 0)) + scale\_fill\_manual(values = c(``coral'',
``skyblue''), guide = ``none'') + theme\_minimal() + labs( title =
``2024 Yılı İllerin NEt Göç Miktarları'', x = ``Net Göç Miktarı (Bin
Kişi)'', y = ``İller'' )

\section{göç nedenleri stacked halde
gösterme}\label{guxf6uxe7-nedenleri-stacked-halde-guxf6sterme}

NetGoc \textless-
read\_excel(``C:/EMU660\_Project/emu660-spring2026-Selindznn/NetGoc.xlsx'',col\_names
= TRUE)

goc\_nedenleri \textless- c(``Tayin/Is Degisikligi'', ``Ise Baslamak/Is
Bulmak'', ``Egitim'', ``Medeni Durum/Ailevi'', ``Daha Iyi Konut/Yasam'',
``Haneye Bagimli Goc'', ``Memlekete Donus'', ``Saglik/Bakim'', ``Ev
Alinmasi'', ``Emeklilik'', ``Diger'', ``Bilinmeyen'')

colnames(NetGoc){[}1:3{]} \textless- c(``Yil'', ``Il'', ``Toplam\_Goc'')
colnames(NetGoc){[}4:15{]} \textless- goc\_nedenleri veri\_2024
\textless- subset(NetGoc, Yil == 2024 \& Il != ``Toplam-Total'')

veri\_uzun \textless- pivot\_longer(veri\_2024, cols = 4:15, names\_to =
``Goc\_Nedeni'', values\_to = ``Miktar'')

veri\_uzun\(Miktar_Bin <- veri_uzun\)Miktar / 1000

goc\_grafigi \textless- ggplot(veri\_uzun, aes(x = Miktar\_Bin, y =
reorder(Il, Toplam\_Goc), fill = Goc\_Nedeni)) + geom\_col(position =
``stack'') + theme\_minimal() + labs( title = ``2024 Yılı İllerin Göç
Nedenlerine Göre Dağılımı'', subtitle = ``Miktarlar bin kişi cinsinden
ifade edilmiştir'', x = ``Net Göç Miktarı (Bin Kişi)'', y = ``İller'',
fill = ``Göç Nedenleri'' ) + theme( legend.position = ``bottom'',
legend.text = element\_text(size = 8), axis.text.y = element\_text(size
= 7) ) + guides(fill = guide\_legend(ncol = 3))

print(goc\_grafigi)

\subsection{3.2 Trend Analysis}\label{trend-analysis}

\#install.packages(``ggrepel'') library(ggplot2) library(readxl)
library(dplyr) library(ggrepel) \# Etiketlerin üst üste binmesini
önlemek için

yasam\_df \textless-
read\_excel(``C:/EMU660\_Project/emu660-spring2026-Selindznn/YasamEndeksi.xlsx'',
col\_names = TRUE) goc\_df \textless-
read\_excel(``C:/EMU660\_Project/emu660-spring2026-Selindznn/Gochizi.xlsx'',
col\_names = TRUE)

colnames(yasam\_df){[}1:2{]} \textless- c(``Il'', ``Endeks'')
colnames(goc\_df){[}1:3{]} \textless- c(``Yil'', ``Il'', ``Goc'')

goc\_2024 \textless- subset(goc\_df, Yil == 2024 \& Il !=
``Toplam-Total'')

veri\_seti \textless- inner\_join(yasam\_df, goc\_2024, by = ``Il'')

grafik \textless- ggplot(veri\_seti, aes(x = Endeks, y = Goc, label =
Il)) + \# label = Il eklendi geom\_point(color = ``\#1f77b4'', size = 3,
alpha = 0.7) + geom\_text\_repel(size = 3, max.overlaps = 15) + \#
Noktaların üzerine il isimlerini akıllıca yerleştirir
geom\_smooth(method = ``lm'', color = ``red'', se = TRUE, formula = y
\textasciitilde{} x) + labs( title = ``İl Yaşam Endeksi ve 2024 Net Göç
Hızı İlişkisi'', x = ``Genel Yaşam Endeksi'', y = ``Net Göç Hızı (‰)'' )
+ theme\_minimal()

print(grafik)

\subsection{3.4 Results}\label{results}

xxxxxx

\section{4. Results and Key Takeaways}\label{results-and-key-takeaways}

xxxxxx




\end{document}
